%==============================================================================
%  24-how-ai-works.tex — How AI Works in Companies
%==============================================================================
\strategyPart{24 \textbar\ How AI Works}{How AI Works in Companies}

AI works in companies by combining \textbf{data, models, software, employees,
processes and governance}. It helps organizations understand information,
predict outcomes, recommend decisions, generate content, automate repetitive
work and execute controlled actions through assistants or agents. It should
not be understood only as a chatbot: in an enterprise it operates at
\textbf{five levels}.

\section{Five levels of enterprise AI}

\begin{itemize}
  \item \textbf{Individual assistance} — helping one employee perform a task.
  \item \textbf{Functional support} — supporting sales, finance, HR or marketing.
  \item \textbf{Workflow automation} — coordinating several steps of an
        end-to-end process.
  \item \textbf{Organizational intelligence} — connecting information and
        decisions across departments.
  \item \textbf{Business-model transformation} — enabling new products,
        services and ways of working.
\end{itemize}

Deloitte's 2026 enterprise research describes the current movement from
experimentation toward \textbf{scaling}, including broader workforce access
and rising interest in autonomous AI agents [Deloitte State of AI in the
Enterprise 2026].

\section{AI as an employee assistant}

\begin{itemize}
  \item \textbf{Knowledge assistant} — searches documents, policies and
        databases; reduces search time, lowers dependency on experts and
        speeds onboarding.
  \item \textbf{Writing \& communication assistant} — drafts emails, reports,
        meeting summaries and proposals; reduces administrative load.
  \item \textbf{Meeting assistant} — transcribes, extracts decisions and
        action items; creates searchable organizational memory.
  \item \textbf{Personal productivity assistant} — organizes tasks,
        prioritizes work and prepares follow-ups.
\end{itemize}

\textbf{Key control:} employees must not treat AI output as automatically
correct; human review remains necessary for important decisions
[HBR 2018; NIST AI RMF].

\section{The central argument}

\begin{center}
  \yBanner[color=inNavy]{AI empowers companies not merely by automating tasks, but by combining assistants, predictive systems, generative models, agents, human expertise and governed data into redesigned organizational workflows.}
\end{center}

This distinction matters because the research shows the greatest value comes
when organizations \textbf{redesign work and operating models}, rather than
simply making existing tasks faster [McKinsey 2025; Deloitte 2026]. McKinsey
reports that nearly two-thirds of surveyed organizations had not yet begun
scaling AI across the enterprise — the redesign is the missing step
[McKinsey State of AI 2025].
